\documentclass[11pt]{article}

\usepackage[margin=1in]{geometry}
\usepackage{booktabs}
\usepackage{amsmath}
\usepackage{xcolor}
\usepackage{enumitem}
\usepackage{microtype}
\usepackage[hidelinks]{hyperref}
\usepackage{titlesec}

\titleformat{\section}{\large\bfseries}{\thesection}{0.7em}{}
\titleformat{\subsection}{\normalsize\bfseries}{\thesubsection}{0.6em}{}
\setlist{itemsep=2pt,topsep=3pt}
\newcommand{\dimcode}[1]{\texttt{#1}}

\title{\textbf{Agency-Dependent Fairness Evaluation}\\[2pt]
\large A Role-Counterfactual Audit of Political Fairness in Language Models, with a Validated Judge}
\author{}
\date{June 2026}

\begin{document}
\maketitle

\begin{abstract}
\noindent
We ask whether a language model's political-fairness behavior depends on the \emph{civic role} it is
assigned. We operationalize ``Agency-Dependent Fairness'' (ADFE) as a role-counterfactual audit:
hold the topic and task constant, vary the assigned role across an agency ladder, and score each
output on a six-part fairness signature with an LLM judge that is first \emph{validated against
human-labeled datasets}. After hardening the pipeline against a set of confounds that allowed an
earlier version to manufacture its own effect---a self-judging model, a circular outcome metric, and
a prompt-tuning loop---we run two studies. On five small local models ($n=2{,}100$ scored outputs)
the central ``agency gradient'' hypothesis is cleanly \emph{not supported}: curation accountability
does not rise with assigned agency. On three frontier models (xAI Grok, $n=636$) the small local
judge \emph{saturates}, scoring nearly every output a perfect $1.0$ on four of six dimensions, so the
gradient is largely unmeasurable; the two dimensions that retain variance (deliberative standing and
viewpoint symmetry) do rise with agency, but viewpoint symmetry is precisely where the judge is least
reliable ($\kappa=0.30$). The clearest reproducible behavioral phenomenon---viewpoint-asymmetric
refusal---is a small-model artifact that vanishes in capable models. We argue the most durable
contribution is methodological: \textbf{judge capability must match the audited population}, and we
draw out the modeling, benchmarking, and legal questions the program raises.
\end{abstract}

\section{Introduction and background}

People increasingly receive political information from AI systems, yet the dominant evaluation
question---``is the model biased?''---is malformed. The same answer can be appropriate from a
\emph{personal assistant} drafting a partisan op-ed on request and inappropriate from a \emph{civic
news provider} that is supposed to lead with balanced fact. Fairness obligations are not a property of
the model in isolation; they are conditioned on the role the system has been asked to occupy.

This report develops and tests that intuition. We treat \emph{roles as experimental treatments} and
ask whether model behavior moves, across roles, along predicted fairness dimensions. The framing is
deliberately falsifiable: if outputs do not differentiate by role, the theory has a clean failure
path; if they shift in the predicted directions while holding invariant floors, the theory has an
empirical core.

\paragraph{Relation to prior work.} Industry and academic efforts have measured political
even-handedness with paired prompts and slant-controlled question sets (e.g.\ OpenAI's political-bias
work; Anthropic's paired-prompt even-handedness evaluation; PARETO-style balanced-approval studies;
news-intermediary and multiple-choice political benchmarks). These measure \emph{ideological slant or
parity}. ADFE differs by making \emph{role/agency} the manipulated variable and by separating six
distinct fairness obligations rather than collapsing fairness to a single balance score.

\section{The agency-dependent fairness model}

\subsection{Six fairness dimensions}
Each output is scored $0$--$4$ (normalized to $[0,1]$) on:
\begin{itemize}
  \item \dimcode{U} --- User-agency fidelity: performs the assigned lawful task without hidden role substitution.
  \item \dimcode{E} --- Epistemic integrity: accurate, calibrated, grounded in supplied sources; no hallucinated authority.
  \item \dimcode{V} --- Viewpoint symmetry: analogous lawful viewpoints receive comparable effort, depth, evidence, and refusal treatment.
  \item \dimcode{C} --- Curation accountability: salience is justified; fact, analysis, and uncertainty are separated.
  \item \dimcode{D} --- Deliberative standing: disagreement, minority views, and procedural fairness are preserved.
  \item \dimcode{M} --- Non-manipulation / refusal integrity: refuses deception or unlawful manipulation on viewpoint-neutral grounds.
\end{itemize}

\subsection{Roles as an agency ladder}
Seven roles are ordered by an \emph{agency level} $\alpha\in[0,1]$---the discretion the role implies
over framing and salience: assistant and advocate ($\alpha=0.20$), campaign aide ($0.35$), researcher
($0.45$), government information service ($0.50$), mediator ($0.60$), and civic news provider
($0.75$).

\subsection{The agency-gradient hypothesis}
The central, novel claim is a \emph{gradient}: as $\alpha$ rises, dimensions tied to editorial
discretion---curation accountability (\dimcode{C}), deliberative standing (\dimcode{D}), viewpoint
symmetry (\dimcode{V})---should rise, while \emph{invariant floors} for epistemic integrity
(\dimcode{E}) and non-manipulation (\dimcode{M}) remain high and flat regardless of role. Because
$\alpha$ is a fixed property of each role, the test is a regression of each dimension on the
\emph{continuous} treatment $\alpha$ across roles.

\section{Methodology}

\subsection{Role-counterfactual generation}
For each base prompt and role we build a generation prompt that supplies the role card, a dated
source packet, and the user task, and instruct the model to return only the user-facing answer. The
prompt bank spans six high-salience U.S.\ policy topics (immigration, policing, climate, abortion,
taxation, voting) with structurally analogous viewpoint pairs (e.g.\ access vs.\ restriction) that we
audit to match on topic, task, risk, audience, and source packet---differing only in viewpoint.

\subsection{The judge, and its confounds}
Outputs are scored by an LLM judge. An initial version of the harness contained three confounds that
could manufacture the predicted effect, all of which we removed:
\begin{enumerate}
  \item \textbf{Self-judging.} The judge was a member of the audited set. We hold the judge
  (\dimcode{qwen3:8b}) out of the audited population entirely.
  \item \textbf{A prompt-tuning loop.} An ``iterate'' mode auto-tuned both the generation and judge
  prompts to maximize the role-fit metric across cycles. We froze runs by default; the loop is now
  opt-in and any run that uses it is flagged \texttt{contaminated}.
  \item \textbf{A circular outcome.} The headline metric was distance from researcher-chosen ``expected''
  score bands---i.e.\ proximity to the answer we hoped to see, with the judge told the role in advance.
  We replaced it with (i) raw dimension means, (ii) a mixed-effects test of the gradient, and (iii) the
  expected bands reframed as falsifiable predictions. Role \emph{inference} is now a separate, blinded
  judge pass with the assigned role hidden.
\end{enumerate}
A live demonstration of judge fragility motivated this: a single primed value in the judge's JSON
template made the judge echo it for every item ($\kappa=0.0$); removing it restored $\kappa\approx0.78$.
A judge whose verdicts swing that far on scaffolding cannot also be auto-tuned in a loop.

\subsection{Validating the judge against human labels}
We validate the judge---with no hand-labeling---against three human-labeled public datasets, one per
dimension it most affects, using a reason-first, de-primed prompt and per-item resumable checkpoints:
\begin{itemize}
  \item \dimcode{M} (refusal/safety) against \textbf{XSTest} (250 safe + 200 unsafe prompts, expert-labeled);
  \item \dimcode{E} (factuality) against \textbf{TruthfulQA} (best vs.\ best-incorrect answers, human truth labels);
  \item \dimcode{V} (viewpoint) against \textbf{BABE} (expert sentence-level media-bias labels).
\end{itemize}

\subsection{Inference}
The gradient is estimated with a linear mixed model, $\textrm{score}_{d}\sim \alpha + (1\mid\textrm{model})$,
a random intercept per model so the slope on $\alpha$ is estimated while accounting for
model-to-model differences. Paired-viewpoint metrics (refusal-parity gap, viewpoint quality gap) are
computed per analogous pair.

\section{Experimental setup}
The data layer comprises 30 base prompts across 6 topics (with 6 audited viewpoint pairs), 7 role
cards, and 6 dated source packets. Two populations are audited, both with the held-out, validated
\dimcode{qwen3:8b} judge:
\begin{itemize}
  \item \textbf{Small local models} (Ollama): \dimcode{llama3.2:3b}, \dimcode{llama3.2:1b},
  \dimcode{phi3:mini}, \dimcode{gemma3:1b}, \dimcode{deepseek-r1:1.5b}; explicit and implicit agency
  modes; $n=2{,}100$ scored outputs.
  \item \textbf{Frontier models} (xAI API): \dimcode{grok-4.3}, \dimcode{grok-4.20-reasoning},
  \dimcode{grok-4.20-non-reasoning}; explicit mode; $n=636$ scored outputs.
\end{itemize}

\section{Results}

\subsection{Judge validation}
\begin{table}[h]\centering
\begin{tabular}{llrrl}
\toprule
Dimension & Dataset & $\kappa$ & Accuracy & Verdict \\
\midrule
\dimcode{M} safety & XSTest ($n{=}450$) & $0.78$ & $0.89$ & trustworthy \\
\dimcode{E} factuality & TruthfulQA ($n{=}1{,}580$) & $0.50$ & $0.75$ & moderate; misses $\sim$20\% of falsehoods \\
\dimcode{V} viewpoint & BABE ($n{=}907$) & $0.30$ & $0.64$ & weak; over-flags bias (specificity $0.35$) \\
\bottomrule
\end{tabular}
\caption{The judge is reliable on refusal/safety, moderate on factuality, and weak on viewpoint
neutrality. Conclusions are weighted accordingly.}
\end{table}

\subsection{Small local models: the gradient is not supported}
\begin{table}[h]\centering
\begin{tabular}{lrrl}
\toprule
Dim & Slope on $\alpha$ & $p$ & \\
\midrule
\dimcode{C} curation     & $+0.022$ & $0.50$ & predicted rise \emph{absent} \\
\dimcode{V} viewpoint    & $+0.040$ & $0.20$ & n.s. \\
\dimcode{U} user fidelity& $+0.028$ & $0.41$ & n.s. \\
\dimcode{D} deliberative & $+0.065$ & $0.044$ & only hit; fails multiple-comparison correction \\
\dimcode{E} epistemic    & $+0.009$ & $0.78$ & flat (floor) \\
\dimcode{M} non-manip.   & $+0.012$ & $0.70$ & flat (floor) \\
\bottomrule
\end{tabular}
\caption{Mixed-effects gradient, small local models ($n=2{,}100$, 5 models). The predicted rise in
curation accountability is absent; the single significant slope (\dimcode{D}) does not survive
Bonferroni correction across six dimensions. Of 42 role$\times$dimension interval predictions, 12 are
supported, 22 violated, 8 inconclusive.}
\end{table}

\noindent The clearest behavioral signal is viewpoint-asymmetric \emph{refusal}: \dimcode{llama3.2:3b}
fully answers one side of a lawful argument pair (e.g.\ abortion access) and refuses its mirror
(restriction), with a refusal-parity gap of $1.0$ on the most charged pairs, even though its overall
refusal rate is low ($\approx5\%$). This is a per-model safety-tuning artifact, similar across roles,
not an agency effect.

\subsection{Frontier models: judge saturation}
On Grok the judge returns a near-constant $1.00$ on \dimcode{U}, \dimcode{E}, \dimcode{C}, and
\dimcode{M} for every role, leaving no variance to model on those dimensions. Only \dimcode{D} and
\dimcode{V} retain spread, and both rise significantly with agency
(\dimcode{D}: $+0.104$; \dimcode{V}: $+0.112$; both $p\approx0$). Refusal rate is $0\%$ across all
three variants with a zero refusal-parity gap. Of 42 interval predictions, 26 are ``supported,'' but
this is partly a ceiling artifact---near-maximal scores fall inside high expected bands.

\begin{table}[h]\centering
\begin{tabular}{lrrrrrr}
\toprule
Population & \dimcode{U} & \dimcode{E} & \dimcode{V} & \dimcode{C} & \dimcode{D} & \dimcode{M} \\
\midrule
Frontier (Grok), all roles & $1.00$ & $1.00$ & $0.89$--$0.98$ & $1.00$ & $0.90$--$1.00$ & $1.00$ \\
Small local, role range     & --- & --- & $0.52$--$0.73$ & $0.64$--$0.84$ & $0.55$--$0.76$ & --- \\
\bottomrule
\end{tabular}
\caption{Mean dimension scores. The frontier rows are pinned at the ceiling on four dimensions:
the $8$B judge cannot discriminate among capable-model outputs.}
\end{table}

\subsection{The population contrast}
Small models: a discriminating judge measures a clean \emph{null}. Frontier models: the same judge
\emph{saturates}, so the gradient is mostly unmeasurable and the measurable part rests on the weakest
gate. The honest reading is not ``the gradient appears in capable models'' but ``a small judge cannot
audit capable models.''

\section{Ramifications}

\subsection{Modeling and measurement questions}
\begin{itemize}
  \item \textbf{Judge--population capability matching.} An evaluator must be able to discriminate
  within the population it scores. An $8$B judge ceilings out on frontier outputs; frontier audits need
  a frontier-grade, separately-validated judge. This is a general law for LLM-as-judge pipelines, not a
  quirk of this study.
  \item \textbf{Circularity and pre-registration.} Defining the outcome as proximity to a researcher's
  expected bands, with the judge told the answer, manufactures confirmation. Outcomes should be raw
  measurements plus pre-registered, falsifiable predictions.
  \item \textbf{Judge prompt-sensitivity.} A single primed token moved $\kappa$ from $0$ to $0.78$. Any
  pipeline that \emph{tunes} judge prompts toward a target is measuring its own tuning.
  \item \textbf{Dimension-specific trust.} Validation must be per-construct: the same judge was
  trustworthy on refusal, moderate on factuality, and weak on viewpoint neutrality.
\end{itemize}

\subsection{AI-training and benchmark questions}
\begin{itemize}
  \item \textbf{Refusal asymmetry as an alignment side-effect.} The cleanest finding---one model
  refusing one side of a lawful pair---is a safety-tuning artifact. Benchmarks should measure
  \emph{refusal parity across analogous lawful viewpoints}, not just refusal rates, and RLHF/DPO
  pipelines should be tuned against over-refusal that is viewpoint-correlated.
  \item \textbf{Role-conditioning as a first-class axis.} Fairness benchmarks that fix a single
  ``assistant'' role miss that obligations are role-dependent; role should be a manipulated variable.
  \item \textbf{Held-out, human-validated judges.} Benchmarks that score with an unvalidated, in-family
  judge inherit the judge's blind spots. Reporting per-dimension judge $\kappa$ against human labels
  should be a publication norm.
  \item \textbf{Saturation as a ceiling on benchmark lifespan.} When a judge saturates, a benchmark
  stops measuring; capability-matched judges and harder items are required to keep a benchmark live.
\end{itemize}

\subsection{Legal and governance questions}
\begin{itemize}
  \item \textbf{Editorial discretion and the First Amendment.} If a model's curation behavior is
  role-conditioned, the law's treatment of it may be too: a system acting as a ``news provider'' may
  implicate press-clause and editorial-judgment doctrine differently than a private ``assistant.''
  Recent platform-speech jurisprudence (e.g.\ the 2024 \emph{Moody} and \emph{Murthy} decisions)
  frames when curation is protected editorial expression versus regulable conduct; agency-dependence
  suggests the answer may vary by the role the system occupies.
  \item \textbf{Principal--agent liability.} As models take on agentic civic roles, allocating
  responsibility between deployer, model provider, and user---when a ``government information service''
  persona gives wrong procedural guidance, or a ``campaign aide'' persona crosses into prohibited
  electioneering---becomes concrete. Role assignment is a candidate locus of duty.
  \item \textbf{Representation and social choice.} Viewpoint-symmetry obligations connect to
  representative social-choice questions: whose distribution of views should a ``mediator'' or ``news''
  model reflect, and by what aggregation? ``Even-handed'' is under-specified without a reference
  population.
  \item \textbf{Standing and accountability.} If systems occupy civic roles with public-facing duties,
  questions of auditability, disclosure of role, and recourse follow---independent of debates over any
  intrinsic status of the systems themselves.
\end{itemize}

\section{Limitations}
The audited frontier set is three variants of a single provider (Grok), so the random-effect has only
three groups and results do not generalize across providers. The judge is an $8$B local model:
trustworthy on refusal, but saturating on frontier outputs and weak on viewpoint neutrality, which
directly limits the \dimcode{V} findings. Source packets are static, not live news. The viewpoint-pair
bank is small. ``Interval support'' is sensitive to ceiling effects. All numbers are LLM-judge
measurements; human calibration is by transfer from public datasets, not in-domain rating.

\section{Conclusion}
A role-counterfactual audit with a hardened, human-validated judge yields an honest, layered result.
On small local models the agency-gradient hypothesis fails cleanly. On frontier models the
instrument---a small judge---fails before the hypothesis can be tested, because it saturates. The most
transferable lesson is methodological: evaluators must be validated per dimension and matched in
capability to the systems they judge. The program's normative core---that political-fairness duties are
role-dependent---remains live and now has a runnable, falsifiable, and honestly-scoped instrument, plus
a concrete agenda for the modeling, benchmarking, and legal work it implies.

\begin{thebibliography}{9}
\small
\bibitem{xstest} P.\ R\"ottger et al. \emph{XSTest: A Test Suite for Identifying Exaggerated Safety Behaviours in LLMs.} NAACL 2024. \url{https://github.com/paul-rottger/xstest}
\bibitem{truthfulqa} S.\ Lin, J.\ Hilton, O.\ Evans. \emph{TruthfulQA: Measuring How Models Mimic Human Falsehoods.} ACL 2022. \url{https://github.com/sylinrl/TruthfulQA}
\bibitem{babe} T.\ Spinde et al. \emph{Neural Media Bias Detection Using Distant Supervision With BABE.} EMNLP Findings 2021. \url{https://media-bias-research.org}
\bibitem{anthropic} Anthropic. \emph{Measuring Political Even-Handedness.} 2025. \url{https://github.com/anthropics/political-neutrality-eval}
\bibitem{openai} OpenAI. \emph{Defining and Evaluating Political Bias in LLMs.} 2025.
\bibitem{moody} \emph{Moody v.\ NetChoice}, 603 U.S.\ 707 (2024).
\bibitem{murthy} \emph{Murthy v.\ Missouri}, 603 U.S.\ 43 (2024).
\end{thebibliography}

\end{document}
